\beginsong{Să stai pe Stânca mântuirii}
\beginverse
Să stai pe Stânca mântuirii
senin, voios și neclintit,
chiar când stihiile naturii
pornite din hotarul firii,
dezlănțuitu-s-au cumplit,
dezlănțuitu-s-au cumplit.
\endverse

\beginverse
Să stai sub adăpost de sânge
statornic, bine-ncredințat.
Când cugetul rănit se frânge
și inima de-amar îți plânge
tu scapi prin Cel crucificat,
tu scapi prin Cel crucificat.
\endverse

\beginverse
Să stai la umbra Celui Veșnic,
odihna Lui să-ți fie scut.
De sunt văpăi, de-i întuneric,
prădat de-ai fi de cel puternic,
să nu te uiți la ce-ai pierdut,
să nu te uiți la ce-ai pierdut.
\endverse

\beginverse
Să stai la post în vremuri grele,
nu dând la bir cu cei ce fug.
Tu, mută-ți cuibul între stele,
lumină dând și tu ca ele,
tu ai viață din belșug,
tu ai viață din belșug.
\endverse

\beginverse
Să stai în strânsă legătură
ca mădular în Trupul Lui.
Și-n dragoste fără măsură
cu-a Duhului Sfânt armătură,
luptând ca fiu al cerului,
luptând ca fiu al cerului!
\endverse

\beginverse
Chiar dacă tremuri în furtună
de val, de vânt vijelios,
de-asupra ta stă-ntinsă-o Mână
de cui străpunsă, dar stăpână!
E Mâna tare-a lui Hristos,
e Mâna tare-a lui Hristos!
\endverse
\endsong